\bookchapter{ภาคผนวก C: ตารางเทียบศัพท์ “ภาษาชาวบ้าน” ฉบับย่อ}

\begin{center}
\begin{tabularx}{\linewidth}{|X|X|X|}
\hline
\textbf{คำที่คุณอาจได้ยิน} & \textbf{แปลว่าจริง ๆ แล้วคือ……} & \textbf{อธิบายในหนึ่งประโยค} \tabularnewline
\hline
SoC / ชิป & สมองและหัวใจของมือถือ & ชิ้นส่วนหลักที่กำหนดว่ามือถือเร็วหรือไม่ เกมกระตุกหรือไม่ \tabularnewline
RAM (หน่วยความจำ) & "โต๊ะทำงานชั่วคราว" ของมือถือ & โต๊ะยิ่งใหญ่ เปิดแอปพร้อมกันได้มากเท่าไร สลับไปมาก็ยิ่งลื่น \tabularnewline
ROM (พื้นที่จัดเก็บ) & "คลังใหญ่" ของมือถือ & ที่เก็บรูป วิดีโอ แอป ตัวเลข 128G, 256G หมายถึงตรงนี้ \tabularnewline
OLED / LCD & สองวัสดุของหน้าจอ & OLED สีสันจัดจ้านกว่าและประหยัดไฟกว่า LCD ส่วน LCD ถนอมตากว่าและต้นทุนต่ำ \tabularnewline
NFC & การสื่อสารไร้สายระยะใกล้ & ฟังก์ชัน "แตะแปะ" ที่ทำให้มือถือของคุณกลายเป็นบัตรโดยสาร บัตรเข้าออก บัตรธนาคารได้ \tabularnewline
การอัปเดต OTA & การอัปเกรดระบบแบบไร้สาย & "การอัปเดตระบบ" ที่กดใน [การตั้งค่า] นั่นแหละคือมัน \tabularnewline
Root & สิทธิ์ผู้ดูแลระบบสูงสุดของ Android & ได้มาแล้วแก้ไขไฟล์ระบบได้ แต่ก็มีความเสี่ยง \tabularnewline
BL Lock & กุญแจ Bootloader & กุญแจความปลอดภัยที่ผู้ผลิตใส่มาให้ ปลดล็อกแล้วจึงจะแฟลชรอมได้ \tabularnewline
NPU & หน่วยประมวลผลโครงข่ายประสาทเทียม & ชิปที่ทำงานเฉพาะด้าน AI ไวกว่า CPU หลายสิบเท่า และกินไฟน้อยกว่า \tabularnewline
AI บนเครื่อง & AI ที่รันในเครื่องมือถือ & AI รันบนมือถือได้ตรง ๆ ไม่ต้องต่อเน็ต เร็วและปกป้องความเป็นส่วนตัวดี \tabularnewline
Qi2 & มาตรฐานการชาร์จไร้สายรุ่นใหม่ & การชาร์จไร้สายที่มีแม่เหล็กช่วยระบุตำแหน่ง คล้าย MagSafe ชาร์จสะดวกขึ้น \tabularnewline
eSIM & ซิมแบบฝังในตัว & ซิมที่ไม่มีการ์ดจริง เขียนลงมือถือตรง ๆ เปลี่ยนผู้ให้บริการโดยไม่ต้องเปลี่ยนการ์ด \tabularnewline
LTPO & เทคโนโลยีปรับอัตรารีเฟรชอัตโนมัติ & อัตรารีเฟรชของหน้าจอปรับตามเนื้อหาอัตโนมัติ (1-120Hz) ลื่นและประหยัดไฟ \tabularnewline
\hline
\end{tabularx}
\end{center}
